% GENERATED FILE. Edit canonical lessons, then run npm run generate:books.
% canonical-source-hash: fnv1a64:aecb8d62ea5a27df
% canonical-lessons: RU-C241-podval, RU-C241-balkon, RU-C241-dvor, RU-C241-naprotiv, RU-C241-mezhdu

\chapter{Basement, Balcony, Yard, Opposite, Between}
\label{ch:ru-basement-balcony-yard-opposite-between}
\hlchaptermodality{241}

\begin{chapteropening}
\textbf{By the end of this chapter:} \emph{I can say a basement, a balcony, a yard, opposite and between.}

Complete the last lesson of chapter 241: I can say a basement, a balcony, a yard, opposite and between.
\end{chapteropening}

\section[\texorpdfstring{\ru{подвал} (podvál)}{подвал (podvál)}]{\ru{подвал} (podvál) — a basement, a cellar}

\label{lesson:RU-C241-podval}



\begin{quote}
Before the new one: say the Russian for at once, then the Russian for usually.
\end{quote}

\subsection*{You'll want to know: \ru{подвал}}
\textbf{\ru{подвал}} (\emph{podvál}) — \textquotedblleft{}a basement, a cellar\textquotedblright{}.

\begin{cousinweb}
\textbf{\ru{подвал}} is \emph{\ru{под}}, \textquotedblleft{}under\textquotedblright{}, with \emph{\ru{вал}}, \textquotedblleft{}a rampart, an earth bank\textquotedblright{}, which Vasmer thinks came through Polish from German \emph{Wall}, Latin \emph{vallum}: a room under the house.
\end{cousinweb}

\begin{grammarlens}[title={What you've built}]
One more word for this chapter.
\end{grammarlens}

\subsection*{Your turn}
\begin{itemize}
\raggedright
  \item \emph{Say it:} \emph{podvál}
  \item \emph{Say it:} \emph{podvál}, once more
  \item \emph{Say it:} say \emph{podvál}
  \item \emph{Recall:} say the Russian for at once, then the Russian for usually, then say \emph{podvál} again
\end{itemize}

\begin{tcolorbox}[breakable,colback=teal!4,colframe=teal!35!black,title={Before you move on}]
What does \ru{подвал} mean? (\textquotedblleft{}A basement, a cellar\textquotedblright{}.) Say it once more.
\end{tcolorbox}
\section[\texorpdfstring{\ru{балкон} (balkón)}{балкон (balkón)}]{\ru{балкон} (balkón) — a balcony}

\label{lesson:RU-C241-balkon}



\begin{quote}
Before the new one: say the Russian for usually, then the Russian for a basement.
\end{quote}

\subsection*{You'll want to know: \ru{балкон}}
\textbf{\ru{балкон}} (\emph{balkón}) — \textquotedblleft{}a balcony\textquotedblright{}.

\begin{cousinweb}
\textbf{\ru{балкон}} — known from 1725; it came from German \emph{Balkon} or French \emph{balcon}, from Italian \emph{balcone}, which goes back to Lombardic and Old High German \emph{balko}, \textquotedblleft{}a beam\textquotedblright{} (Vasmer).
\end{cousinweb}

\begin{grammarlens}[title={What you've built}]
One more word for this chapter.
\end{grammarlens}

\subsection*{Your turn}
\begin{itemize}
\raggedright
  \item \emph{Say it:} \emph{balkón}
  \item \emph{Say it:} \emph{balkón}, once more
  \item \emph{Say it:} say \emph{balkón}
  \item \emph{Recall:} say the Russian for usually, then the Russian for a basement, then say \emph{balkón} again
\end{itemize}

\begin{tcolorbox}[breakable,colback=teal!4,colframe=teal!35!black,title={Before you move on}]
What does \ru{балкон} mean? (\textquotedblleft{}A balcony\textquotedblright{}.) Say it once more.
\end{tcolorbox}
\section[\texorpdfstring{\ru{двор} (dvor)}{двор (dvor)}]{\ru{двор} (dvor) — a yard, a courtyard}

\label{lesson:RU-C241-dvor}



\begin{quote}
Before the new one: say the Russian for a basement, then the Russian for a balcony.
\end{quote}

\subsection*{You'll want to know: \ru{двор}}
\textbf{\ru{двор}} (\emph{dvor}) — \textquotedblleft{}a yard, a courtyard\textquotedblright{}.

\begin{cousinweb}
\textbf{\ru{двор}} $\leftarrow$ Proto-Slavic \emph{*dvorŭ} — inherited, with \emph{\ru{дверь}}, \textquotedblleft{}door\textquotedblright{} (Vasmer). Vasmer relates it to Latvian \emph{dvars}, \textquotedblleft{}a gate\textquotedblright{}, and Sanskrit \emph{dvāra-}, \textquotedblleft{}gate\textquotedblright{}: a yard is the space inside the gate.
\end{cousinweb}

\begin{grammarlens}[title={What you've built}]
One more word for this chapter.
\end{grammarlens}

\subsection*{Your turn}
\begin{itemize}
\raggedright
  \item \emph{Say it:} \emph{dvor}
  \item \emph{Say it:} \emph{dvor}, once more
  \item \emph{Say it:} say \emph{dvor}
  \item \emph{Recall:} say the Russian for a basement, then the Russian for a balcony, then say \emph{dvor} again
\end{itemize}

\begin{tcolorbox}[breakable,colback=teal!4,colframe=teal!35!black,title={Before you move on}]
What does \ru{двор} mean? (\textquotedblleft{}A yard, a courtyard\textquotedblright{}.) Say it once more.
\end{tcolorbox}
\section[\texorpdfstring{\ru{напротив} (naprótiv)}{напротив (naprótiv)}]{\ru{напротив} (naprótiv) — opposite}

\label{lesson:RU-C241-naprotiv}



\begin{quote}
Before the new one: say the Russian for a balcony, then the Russian for a yard.
\end{quote}

\subsection*{You'll want to know: \ru{напротив}}
\textbf{\ru{напротив}} (\emph{naprótiv}) — \textquotedblleft{}opposite\textquotedblright{}.

\begin{cousinweb}
\textbf{\ru{напротив}} is \emph{\ru{на}}, \textquotedblleft{}on\textquotedblright{}, with \emph{\ru{против}}, \textquotedblleft{}against\textquotedblright{}, Proto-Slavic \emph{*protivŭ}, from \emph{*proti}, which Vasmer relates to Latvian \emph{pretī}, \textquotedblleft{}towards, opposite\textquotedblright{}.
\end{cousinweb}

\begin{grammarlens}[title={What you've built}]
One more word for this chapter.
\end{grammarlens}

\subsection*{Your turn}
\begin{itemize}
\raggedright
  \item \emph{Say it:} \emph{naprótiv}
  \item \emph{Say it:} \emph{naprótiv}, once more
  \item \emph{Say it:} say \emph{naprótiv}
  \item \emph{Recall:} say the Russian for a balcony, then the Russian for a yard, then say \emph{naprótiv} again
\end{itemize}

\begin{tcolorbox}[breakable,colback=teal!4,colframe=teal!35!black,title={Before you move on}]
What does \ru{напротив} mean? (\textquotedblleft{}Opposite\textquotedblright{}.) Say it once more.
\end{tcolorbox}
\section[\texorpdfstring{\ru{между} (mézhdu)}{между (mézhdu)}]{\ru{между} (mézhdu) — between}

\label{lesson:RU-C241-mezhdu}



\begin{quote}
Before the new one: say the Russian for a yard, then the Russian for opposite.
\end{quote}

\subsection*{You'll want to know: \ru{между}}
\textbf{\ru{между}} (\emph{mézhdu}) — \textquotedblleft{}between\textquotedblright{}.

\begin{cousinweb}
\textbf{\ru{между}} — taken from Church Slavonic, says Vasmer, beside native \emph{\ru{межу}}, from \emph{\ru{межа}}, \textquotedblleft{}a boundary strip\textquotedblright{}. Vasmer relates \emph{\ru{межа}} to Sanskrit \emph{mádhya-}, Greek \emph{mésos} and Latin \emph{medius}, \textquotedblleft{}middle\textquotedblright{}.
\end{cousinweb}

\begin{grammarlens}[title={What you've built}]
One more word for this chapter.
\end{grammarlens}

\subsection*{Your turn}
\begin{itemize}
\raggedright
  \item \emph{Say it:} \emph{mézhdu}
  \item \emph{Say it:} \emph{mézhdu}, once more
  \item \emph{Say it:} say \emph{mézhdu}
  \item \emph{Recall:} say the Russian for a yard, then the Russian for opposite, then say \emph{mézhdu} again
\end{itemize}

\begin{tcolorbox}[breakable,colback=teal!4,colframe=teal!35!black,title={Before you move on}]
What does \ru{между} mean? (\textquotedblleft{}Between\textquotedblright{}.) Say it once more.
\end{tcolorbox}
